\section{Musculoskeletal X-Ray}


\subsubsection*{What Is This Test?}
Musculoskeletal X-rays image bones and joints to detect fractures, dislocations, arthritis, infections, and tumors. They're the first-line imaging for bone and joint problems.

\subsubsection*{Alternative Names}
MSK X-ray, bone X-ray, extremity X-ray, skeletal radiograph, orthopedic X-ray


\subsubsection*{Principle}
Musculoskeletal X-ray uses radiographic imaging with at least two orthogonal views (AP and lateral) of the bone or joint. The X-ray beam is attenuated differently by cortical bone, trabecular bone, soft tissue, and fat. Fractures appear as radiolucent lines interrupting the cortical bone continuity.

\subsubsection*{History}
X-ray imaging of bones and joints was the first clinical application of radiology, beginning in 1896. It remains the primary imaging modality for fracture detection and joint assessment.

\subsubsection*{Why This Test Is Ordered}
Your doctor orders an X-ray of a bone or joint when you have pain, swelling, or an injury to that area. It answers the question: do you have a fracture, arthritis, dislocation, or bone tumor?

\subsubsection*{Is This a Primary or Secondary Test}
A primary test. X-ray is the first imaging test for bone and joint problems.

\subsubsection*{Limits of musculoskeletal radiographs}
Radiographs show many displaced fractures and established joint changes but can miss nondisplaced or early injuries. Accuracy depends on anatomy, projections, and timing; persistent focal symptoms may need additional imaging.

\subsubsection*{How This Test Is Performed}
The affected body part is positioned against the X-ray detector. Usually two views are taken at different angles (front and side) to provide a complete picture. The test takes about ten to fifteen minutes.

\subsubsection*{What Preparation Is Needed}
Remove any jewelry or metal objects near the area being X-rayed. No other preparation is needed.

\subsubsection*{Contraindications and Precautions}
Pregnancy is not an automatic contraindication; the examination should be justified and optimized, and alternatives considered when appropriate. For contrast-enhanced studies, review the indication, prior severe reactions, and renal function according to local protocol.

\subsubsection*{Risks}
Risks associated with X-ray imaging include:
\begin{itemize}
    \item \textbf{Radiation exposure}: Exposure to low levels of ionizing radiation, with cumulative lifetime dose optimization guided by the ALARA (As Low As Reasonably Achievable) principle.
    \item \textbf{Contrast reaction}: Allergic reactions can occur, including rare severe reactions. Kidney injury after iodinated contrast is uncommon and risk depends on baseline kidney function, illness, and the clinical context.
\end{itemize}


Results are typically available within thirty minutes. You may be placed in a splint or cast while waiting for results if a fracture is suspected.

\subsubsection*{Understanding Your Results}

\subsubsection*{Normal Results}
A normal result shows normal bone alignment, no fractures, no joint abnormalities, and normal soft tissues.
\clearpage
